\documentclass[12pt]{article}
\usepackage[utf8]{inputenc}
\usepackage[T1]{fontenc}
\usepackage{hyperref}
\title{Udihe Language Resources for K-12 Curriculum Development}
\author{Compiled by Claude Sonnet 5 (claude-sonnet-5) for Andr\'as Kornai}
\date{}
\begin{document}
\maketitle

\section{Introduction}
Udihe (Udege, Udeghe; ISO 639-3 ude; Glottocode udih1248) is a Tungusic language of the Russian Far East. Udihe villages are in Khabarovsk Territory (Gvasyugi in the Lazo district, Arsenyevo/Rassvet in the Nanai district) and in Primorsky Territory (Krasny Yar in the Pozharsky district, Akzu in the Terney district). The 2020 census figures are reported as inflated: 679 people gave Udege as their native language, but only about 193 actually speak it. The shift to Russian has been faster than almost anywhere else in Siberia.

\section{Primary Educational Resources}
A Russian Academy of Sciences page says that manuals on Udihe for educational institutions have been published, based on one of two competing alphabets. It names no titles. Wikipedia's Udege alphabets page shows the first page of a handwritten primer by Jhansi Kimonko, but gives no publication details.
\begin{itemize}
\item Udege alphabets (history of the writing system): \url{https://en.wikipedia.org/wiki/Udege_alphabets}
\item The Udeghe Language (Institute of Linguistics, Russian Academy of Sciences): \url{https://lingsib.iea.ras.ru/en/languages/udeghe.shtml}
\end{itemize}

\section{Linguistic and Reference Materials}
Writing began in Latin script in 1931--1937. Since the late 1980s it has used Cyrillic, in two versions for two emerging literary languages, with no generally accepted norm. No dictionary was named in the search results.
\begin{itemize}
\item Udihe bibliography (Surrey Siberian Languages): \url{https://siberianlanguages.surrey.ac.uk/resource/udihe-bibliography/}
\item Atlas of Languages of the Peoples of Russia, Udege: \url{https://atlaskmns.ru/page/en/lang_udegeicy_all.html}
\end{itemize}

\section{Religious and Community Resources}
\begin{itemize}
\item The Red Book of the Peoples of the Russian Empire, Udeghes: \url{https://www.eki.ee/books/redbook/udeghes.shtml}
\end{itemize}

\section{Digital and Audio Resources}
\begin{itemize}
\item Endangered Languages Project entry on Udihe: \url{https://www.endangeredlanguages.com/lang/3021}
\end{itemize}
No audio or software was confirmed in the searches.

\section{Curriculum Development Context}
Only a few dozen to about 200 speakers remain, and most are older adults. A curriculum would therefore be a revitalization programme and would need an agreed alphabet. The first steps are to locate the published school manuals and any dictionary, using Russian-language catalogues, and to document the speakers who remain.
\end{document}